Verification can improve an agent's result, consume resources without benefit,
or remove a more valuable opportunity to check later.  We study a fixed
verification-and-response regime with at most one use across two online
checkpoints.  A paired branch-tree design compares verifying now, reserving
verification for later, and never verifying on the original root population,
including early termination and infeasibility.  Policy assessment separates
sampling uncertainty from an explicitly conditional allowance for
outcome-relevant reconstruction error.  The planned evaluation tests whether
preserving a later verification opportunity improves bounded task utility
against a matched repeated-myopic controller and a context-dependent VOI
scheduler at the same allowed
verifier and total-work budgets, and whether conclusions survive plausible
reconstruction error.  No live benefit or safety result is claimed before this
evaluation: every quantitative statement in this draft is \emph{planned}.
